%!TeX root=../main.tex

\chapter{Introduzione}\label{ch:introduzione}

\section{Motivazioni e contesto}\label{sec:motivazioni-e-contesto}
La gestione di \emph{contest} in tempo reale rappresenta a oggi un processo frammentato e complesso.
Organizzatori, partecipanti e giurati spesso ricorrono a strumenti diversi: servizi di posta elettronica per la comunicazione degli inviti, piattaforme di file sharing per la raccolta dei lavori, fogli di calcolo per la tabulazione dei punteggi.

Questa frammentazione genera inefficienze operative, rischi di errore umano e una significativa difficoltà nel garantire trasparenza e sicurezza.
L'analisi dello stato dell'arte, come verrà dettagliato nel capitolo successivo, ha inoltre rivelato un \textbf{significativo vuoto di mercato}: non esistono soluzioni verticali e accessibili che coprano in modo integrato l'intero ciclo di vita di un \emph{contest}.

Il progetto \textbf{Swift Contest} nasce proprio per colmare questa lacuna, proponendo un sistema moderno, sicuro e multipiattaforma che unisce in un'unica applicazione le principali fasi del ciclo di vita di un \emph{contest}: dalla creazione e configurazione, all’invito dei partecipanti e dei giurati, fino alla raccolta dei lavori e alla gestione delle sessioni di voto in tempo reale.

\section{Obiettivi del progetto}\label{sec:obiettivi-del-progetto}
L’obiettivo primario di questo lavoro è stato quello di \textbf{progettare e implementare l'architettura} di un sistema completo di \emph{contest management}, adottando tecnologie moderne e approcci architetturali rigorosi.

Gli obiettivi specifici, che hanno guidato la progettazione, possono essere così sintetizzati:
\begin{itemize}
    \item \textbf{Sviluppo cross-platform}: garantire un'esperienza utente fluida e unificata su più piattaforme (mobile e web) attraverso lo sviluppo di un client basato su un'unica codebase, un obiettivo reso possibile dall'adozione del framework Flutter\cite{flutter_docs}.
    \item \textbf{Sicurezza ``\emph{by design}''}: adottare un'architettura di sicurezza basata sul \textbf{Principio del Minimo Privilegio (Principle of Least Privilege)}\cite{saltzer1975protection}, dove l'accesso al database è negato per default e concesso selettivamente solo tramite un'API controllata.
    \item \textbf{Esperienza utente mirata}: progettare interfacce utente intuitive e funzionali, modellate sui ruoli specifici dei diversi attori del sistema (organizzatore, partecipante, giurato).
    \item \textbf{Accesso flessibile alla votazione}: implementare un sistema di voto che supporti due modalità di accesso per i giurati.
    Oltre ai giurati nominati (utenti registrati e invitati), il sistema deve permettere la partecipazione di \textbf{giurati semplici} tramite un token di accesso sicuro, abbassando così la barriera d'ingresso per i giurati occasionali.
    \item \textbf{Affidabilità e auditabilità del processo}: assicurare l'integrità e la verificabilità del processo di voto attraverso meccanismi tecnici robusti, come lo \emph{snapshot} dei dati al momento dell'avvio della sessione e la possibilità di applicare restrizioni basate sulla geolocalizzazione.
\end{itemize}

\section{Convenzioni linguistiche}\label{sec:convenzioni-linguistiche}
Al fine di garantire la massima chiarezza e coerenza con gli standard del settore, nel presente elaborato è stata adottata una convenzione linguistica specifica per gli elementi prettamente tecnici.

Per tutti gli artefatti di progettazione e implementazione, come i \textbf{diagrammi UML}, le \textbf{figure architetturali}, i \textbf{listati di codice} e i nomi delle componenti software, è stato scelto di utilizzare la \textbf{lingua inglese}.
Questa decisione riflette lo standard de facto della comunità informatica internazionale, assicurando che la documentazione tecnica sia immediatamente comprensibile e allineata alle convenzioni del settore.
Nel testo descrittivo, termini italiani e inglesi potranno essere usati come sinonimi dove appropriato (es.\ contest/concorso).

\section{Contributi progettuali e ingegneristici}\label{sec:contributi-progettuali-e-ingegneristici}
L'obiettivo primario di questo lavoro di tesi non è di natura teorica, ma \textbf{pratica e orientata al prodotto}: la creazione di una soluzione completa e funzionante per un problema reale e non indirizzato dal mercato attuale.
Il valore del progetto risiede quindi nella dimostrazione di come, partendo da un'esigenza concreta, sia stato possibile progettare e realizzare un sistema software complesso, utilizzando principi di ingegneria moderni come mezzo per garantire la qualità e la robustezza del risultato finale.

Il progetto \textbf{Swift Contest} si presenta quindi come un \emph{blueprint}, ovvero un modello di riferimento, per la creazione di applicazioni collaborative complesse e sicure.
I contributi specifici possono essere così riassunti:
\begin{itemize}
    \item \textbf{Proposta e validazione di una nuova categoria di prodotto}: data l'assenza di soluzioni verticali, il progetto stesso si configura come un caso di studio e una validazione pratica per una nuova categoria di software dedicato alla gestione dei \emph{contest}.
    \item \textbf{Dimostrazione di un'architettura sicura su BaaS}: il sistema dimostra come implementare un'architettura basata sul Principio del Minimo Privilegio su una piattaforma \textbf{Backend-as-a-Service (BaaS)} come Supabase\cite{supabase_docs}, ma \textbf{centralizzando la logica di autorizzazione in API serverless (RPC ed Edge Functions) invece di affidarsi unicamente a regole a livello di database (RLS)}.
    \item \textbf{Modello dati robusto per processi immutabili}: viene presentato un design del database che garantisce l'integrità e l'attendibilità dei processi di voto tramite l'uso di \emph{attributi di snapshot} e denormalizzazione controllata.
\end{itemize}
In sintesi, il valore della tesi è nel dimostrare l'intero processo ingegneristico, dall'analisi di un vuoto di mercato alla progettazione e realizzazione di una soluzione \emph{completa}, \emph{sicura} e \emph{ben ingegnerizzata}.

\section{Struttura del documento}\label{sec:struttura-del-documento}
Il presente elaborato è organizzato in sei capitoli principali, che guidano il lettore attraverso le fasi di analisi, selezione tecnologica, progettazione, implementazione e valutazione del progetto.
\begin{itemize}
    \item \textbf{Capitolo 1 – Introduzione}: illustra le motivazioni, gli obiettivi di alto livello e il valore del lavoro di tesi.
    \item \textbf{Capitolo 2 – Contesto e analisi del sistema}: definisce in dettaglio il problema, analizza lo stato dell’arte e presenta i requisiti del sistema.
    \item \textbf{Capitolo 3 – Tecnologie e strumenti utilizzati}: descrive lo stack tecnologico adottato, giustificando le scelte chiave in relazione agli obiettivi del progetto e ai vincoli di sviluppo.
    \item \textbf{Capitolo 4 – Progettazione e implementazione del sistema}: descrive l’architettura software, il design pattern adottato, la struttura del database e l’implementazione delle principali funzionalità, con esempi di codice e diagrammi tecnici.
    \item \textbf{Capitolo 5 – Distribuzione e prodotto finale}: presenta il processo di distribuzione, l’interfaccia utente e i flussi di utilizzo reali.
    \item \textbf{Capitolo 6 – Conclusioni}: riflette sui risultati ottenuti, ne analizza i limiti e propone possibili sviluppi futuri.
\end{itemize}